\begin{frame}{자주 하는 실수: 공유 공분산 가정}
  GDA는 두 클래스가 \textbf{같은 모양}(공분산)으로 퍼져 있다고
  가정 --- 실제로는 한 클래스가 훨씬 넓게 퍼져 있으면 깨진다.
  \vskip0.4em
  \begin{itemize}
    \item 예: 불합격 $[2,3,4]$ (촘촘) vs 합격 $[6.5,7,10.5]$
      (편차 큼) 이면 각각
      $\sigma_0^2\approx0.667$ vs $\sigma_1^2\approx3.167$
      --- \textbf{약 5배 차이}.
    \item 공유 공분산을 강제하면, 넓게 퍼진 클래스의 ``진짜 모양''을
      무시하고 \textbf{평균적 모양 하나로 뭉뚱그림} $\to$
      결정 경계가 부정확.
    \item \kb{QDA}(이차판별분석): 클래스마다 \textbf{다른} 공분산
      $\Sigma_0\ne\Sigma_1$ 허용 $\to$ 경계가 \textbf{곡선}(이차식).
    \item 트레이드오프: 추정 파라미터가 \textbf{늘어나} 더 많은
      데이터가 필요.
  \end{itemize}
\end{frame}
